\scanpage{60}
\deleted{inverse de la section $b'_0$ de $(\frac12 L)/L$ dont l'image par l'isom can. (induit par $2\,\mathrm{id}_{\frac12 L}$) $(\frac12 L)/L\longrightarrow L/2L$ est $b_0$, \deleted{Alors le torseur} \inserted{\deleted{$\underline L$}} \deleted{$2F$ est le sous faisceau de $\underline L$ engendré par et on a}}

Pour a), \unsure{OPS|DPS} $S$ affine, donc les \inserted{rev. quadratiques} $\Delta$ \deleted{envisagés} de $S$ correspondent à des $\underline L$ torseurs $F$ triviaux, donc $\Delta$ \inserted{isom à un schéma} de la forme
\[\Delta(b,c)\subset\underline L,\quad\Delta(b,c)=\{x\mid x^2-bx+c=0\}\]
où $b\in\Gamma(S,\underline L)$, $c\in\Gamma(S,\underline L^{\otimes2})$. Considérons (\unsure{pour|donc} $\underline L,D,b_0$ fixés) la cat. des $\Delta$ correspondants, \deleted{\unsure{Il}} équivalente à celle des $\Delta(b,c)$, avec $b|S_0=b_0$, $b^2-4c=D$, il faut montrer que cette catégorie, \deleted{\unsure{si}} \inserted{qnd} elle est non vide, est équiv. à la cat. ponctuelle, i.e. \inserted{entre} \deleted{\unsure{deux}} de ces objets il existe un isom. unique. Or les \deleted{\unsure{obj|dits}} isom \inserted{entre $\Delta(b,c)$ et $\Delta(b',c')$} correspondent aux $\tau\in\Gamma(S,\underline L)$ tels que
\[\left\{\begin{aligned} b'&=b+2\tau\\ c'&=c+\tau b+\tau^2\end{aligned}\right.\]

Mais \deleted{$b'\equiv b\ (2)$ \unsure{i.e. que|i.e.}} \inserted{$b'_{S_0}=b_{S_0}\;(=b_0)$} signifie bien que $b'=b+2\tau$ (où $\tau$ est, du reste, déterminé de façon unique, 2 étant $\underline O_S$-régulier). \deleted{Donc, quitte à} \deleted{On a alors}

Considérons $c'_1=c+\tau b+\tau^{2}$, on \deleted{\unsure{a}} donc
\[b^2-4c={b'}^2-4{c'_1}^{\deleted{\unsure{2}}}=D\]
et comme $b'^{2}-4c'=D$, on trouve donc $4(c'_1-c')=0$
\[c'_1-c'=0\qquad\text{i.e.}\quad c'\overset{\unsure{d}}=c',\]
% archive sheet number: 59
